% --- Archon physics formalization source begin ---
% archon:physics
% archon:covers PhyXMiniProblems/problem_phyx_mini_0608.lean
% archon:source-report reports/phyx_mini/problem_phyx_mini_0608.source.json

\chapter{Physics problem phyx\_mini\_0608}
\label{ch:PhyXMiniProblems_problem_phyx_mini_0608}

\paragraph{Problem source.}
\#\# Physical scenario

The starships of the Solar Federation are marked with the symbol of the federation, a circle, while starships of the Denebian Empire are marked with the empire’s symbol, an ellipse whose major axis is 1.40 times longer than its minor axis (\$a = 1.40b\$ in figure).Its marking must be mistaken for that of a federation ship.

\#\# Question

How fast, relative to an observer, does an empire ship have to travel?

\#\# Answer choices

- A: A: \$2.10 \textbackslash{}times 10\textasciicircum{}\{8\}\$ m/s

- B: B: \$2.12 \textbackslash{}times 10\textasciicircum{}\{7\}\$ m/s

- C: C: \$2.13 \textbackslash{}times 10\textasciicircum{}\{6\}\$ m/s

- D: D: \$2.11 \textbackslash{}times 10\textasciicircum{}\{9\}\$ m/s

\#\# Figure caption (auxiliary; use the image as primary evidence)

The image features two stylized spaceship silhouettes, each labeled with corresponding names underneath:

1. **Federation**: 
   - The silhouette is grey and shaped like a delta with smooth edges.
   - It includes a large central circle outlined in purple.
   - The overall design is symmetric, with a pointed top.

2. **Empire**:
   - This silhouette is beige with a more angular shape, featuring spikes at the top.
   - It also has a large central circle outlined in purple.
   - Two dimensions are labeled with arrows:
     - **"a"** indicates a vertical distance from the top of the circle to the top of the shape.
     - **"b"** shows a horizontal measurement at the base of the shape.

Both figures are abstract representations, suggesting different stylistic designs or symbols associated with "Federation" and "Empire."

\#\# Dataset metadata

Relativity; Physical Model Grounding Reasoning, Numerical Reasoning

\paragraph{Recorded answer/context.}
A: A: \$2.10 \textbackslash{}times 10\textasciicircum{}\{8\}\$ m/s

\paragraph{Figure/image path.}
/root/proposal\_for\_physic/science-mango/phyx\_mini\_run/phyx\_data/test\_image/608.png

\paragraph{Formalization target.}
create a compiling Lean file with sorry bodies at `PhyXMiniProblems/problem\_phyx\_mini\_0608.lean`. The Lean declarations must preserve the physical quantities, dimensions or dimensional roles, figure labels, governing-law hypotheses, and final relation expressed by this problem.
Use Mathlib/Physlib names found through LeanExplore where available. If a physics API is missing, introduce faithful local abstractions rather than scalar placeholder aliases.

\begin{theorem}[Physics formalization target]
\label{thm:physics:phyx_mini_0608:target}
\lean{PhyXMiniProblems.ProblemPhyXMini0608.problem_phyx_mini_0608}
\uses{def:physics:phyx-mini-0608:phyxminiproblems-problemphyxmini0608-empiremarkingobservation, def:physics:phyx-mini-0608:phyxminiproblems-problemphyxmini0608-speedfractionoflight, def:physics:phyx-mini-0608:phyxminiproblems-problemphyxmini0608-matchesempiremarkingscenario, def:physics:phyx-mini-0608:phyxminiproblems-problemphyxmini0608-matchessuppliedstarshipfigure, def:physics:phyx-mini-0608:phyxminiproblems-problemphyxmini0608-matchesproblemellipseratio, def:physics:phyx-mini-0608:phyxminiproblems-problemphyxmini0608-hasphysicalempiremarkingparameters, def:physics:phyx-mini-0608:phyxminiproblems-problemphyxmini0608-appearsasfederationcircle, def:physics:phyx-mini-0608:phyxminiproblems-problemphyxmini0608-satisfiesrelativisticaxiscontraction, def:physics:phyx-mini-0608:phyxminiproblems-problemphyxmini0608-matchesdisplayedspeedchoice}
The assigned autoformalize agent should translate this physics problem into Lean declarations in the covered file, with theorem and lemma proof bodies written as `by sorry`.
\end{theorem}
\begin{proof}
This is an autoformalization task, not a proof task. Produce statements that can later be proved by the physics prover without weakening the physical model.
\end{proof}

% --- Archon declaration topology begin ---
\section{Declaration topology}
\begin{definition}[Length Quantity]
  \label{def:physics:phyx-mini-0608:phyxminiproblems-problemphyxmini0608-lengthquantity}
  \lean{PhyXMiniProblems.ProblemPhyXMini0608.LengthQuantity}
  A nonnegative physical length, independent of a choice of units
\end{definition}
\begin{proof}
  This is the declared model definition of Length Quantity.
\end{proof}
\begin{definition}[Speed Quantity]
  \label{def:physics:phyx-mini-0608:phyxminiproblems-problemphyxmini0608-speedquantity}
  \lean{PhyXMiniProblems.ProblemPhyXMini0608.SpeedQuantity}
  A nonnegative, dimensionful speed magnitude
\end{definition}
\begin{proof}
  This is the declared model definition of Speed Quantity.
\end{proof}
\begin{definition}[length Readout]
  \label{def:physics:phyx-mini-0608:phyxminiproblems-problemphyxmini0608-lengthreadout}
  \lean{PhyXMiniProblems.ProblemPhyXMini0608.lengthReadout}
  \uses{def:physics:phyx-mini-0608:phyxminiproblems-problemphyxmini0608-lengthquantity}
  Read a physical length as a real number in a selected length unit
\end{definition}
\begin{proof}
  This is the declared model definition of length Readout.
\end{proof}
\begin{definition}[speed In Meters Per Second]
  \label{def:physics:phyx-mini-0608:phyxminiproblems-problemphyxmini0608-speedinmeterspersecond}
  \lean{PhyXMiniProblems.ProblemPhyXMini0608.speedInMetersPerSecond}
  \uses{def:physics:phyx-mini-0608:phyxminiproblems-problemphyxmini0608-speedquantity}
  Read a physical speed in metres per second
\end{definition}
\begin{proof}
  This is the declared model definition of speed In Meters Per Second.
\end{proof}
\begin{definition}[vacuum Speed Of Light In Meters Per Second]
  \label{def:physics:phyx-mini-0608:phyxminiproblems-problemphyxmini0608-vacuumspeedoflightinmeterspersecond}
  \lean{PhyXMiniProblems.ProblemPhyXMini0608.vacuumSpeedOfLightInMetersPerSecond}
  Physlib's exact vacuum speed of light, read in metres per second
\end{definition}
\begin{proof}
  This is the declared model definition of vacuum Speed Of Light In Meters Per Second.
\end{proof}
\begin{definition}[Starship Affiliation]
  \label{def:physics:phyx-mini-0608:phyxminiproblems-problemphyxmini0608-starshipaffiliation}
  \lean{PhyXMiniProblems.ProblemPhyXMini0608.StarshipAffiliation}
  The two organizations identified by captions in the supplied image
\end{definition}
\begin{proof}
  This is the declared model definition of Starship Affiliation.
\end{proof}
\begin{definition}[Figure Side]
  \label{def:physics:phyx-mini-0608:phyxminiproblems-problemphyxmini0608-figureside}
  \lean{PhyXMiniProblems.ProblemPhyXMini0608.FigureSide}
  The two sides occupied by the ships in the supplied image
\end{definition}
\begin{proof}
  This is the declared model definition of Figure Side.
\end{proof}
\begin{definition}[Marking Shape]
  \label{def:physics:phyx-mini-0608:phyxminiproblems-problemphyxmini0608-markingshape}
  \lean{PhyXMiniProblems.ProblemPhyXMini0608.MarkingShape}
  Geometric kinds of the two affiliation markings
\end{definition}
\begin{proof}
  This is the declared model definition of Marking Shape.
\end{proof}
\begin{definition}[Figure Axis Label]
  \label{def:physics:phyx-mini-0608:phyxminiproblems-problemphyxmini0608-figureaxislabel}
  \lean{PhyXMiniProblems.ProblemPhyXMini0608.FigureAxisLabel}
  The two axis labels printed next to the Empire ellipse
\end{definition}
\begin{proof}
  This is the declared model definition of Figure Axis Label.
\end{proof}
\begin{definition}[Figure Orientation]
  \label{def:physics:phyx-mini-0608:phyxminiproblems-problemphyxmini0608-figureorientation}
  \lean{PhyXMiniProblems.ProblemPhyXMini0608.FigureOrientation}
  Orientation of a labelled axis in the supplied raster
\end{definition}
\begin{proof}
  This is the declared model definition of Figure Orientation.
\end{proof}
\begin{definition}[Ellipse Axis Role]
  \label{def:physics:phyx-mini-0608:phyxminiproblems-problemphyxmini0608-ellipseaxisrole}
  \lean{PhyXMiniProblems.ProblemPhyXMini0608.EllipseAxisRole}
  Physical role assigned to each labelled ellipse axis
\end{definition}
\begin{proof}
  This is the declared model definition of Ellipse Axis Role.
\end{proof}
\begin{definition}[Inertial Frame Label]
  \label{def:physics:phyx-mini-0608:phyxminiproblems-problemphyxmini0608-inertialframelabel}
  \lean{PhyXMiniProblems.ProblemPhyXMini0608.InertialFrameLabel}
  The ship rest frame and the inertial frame of the observing viewer
\end{definition}
\begin{proof}
  This is the declared model definition of Inertial Frame Label.
\end{proof}
\begin{definition}[Starship Marking Figure]
  \label{def:physics:phyx-mini-0608:phyxminiproblems-problemphyxmini0608-starshipmarkingfigure}
  \lean{PhyXMiniProblems.ProblemPhyXMini0608.StarshipMarkingFigure}
  \uses{def:physics:phyx-mini-0608:phyxminiproblems-problemphyxmini0608-starshipaffiliation, def:physics:phyx-mini-0608:phyxminiproblems-problemphyxmini0608-figureside, def:physics:phyx-mini-0608:phyxminiproblems-problemphyxmini0608-markingshape, def:physics:phyx-mini-0608:phyxminiproblems-problemphyxmini0608-figureaxislabel, def:physics:phyx-mini-0608:phyxminiproblems-problemphyxmini0608-figureorientation, def:physics:phyx-mini-0608:phyxminiproblems-problemphyxmini0608-ellipseaxisrole}
  Typed transcription of the primary image. It records the left/right affiliations, the circle/ellipse distinction, and the a and b labels and orientations. The prose-supplied numerical ratio is kept out of this image record because it is not printed in the raster itself
\end{definition}
\begin{proof}
  This is the declared model definition of Starship Marking Figure.
\end{proof}
\begin{definition}[Empire Marking Observation]
  \label{def:physics:phyx-mini-0608:phyxminiproblems-problemphyxmini0608-empiremarkingobservation}
  \lean{PhyXMiniProblems.ProblemPhyXMini0608.EmpireMarkingObservation}
  \uses{def:physics:phyx-mini-0608:phyxminiproblems-problemphyxmini0608-lengthquantity, def:physics:phyx-mini-0608:phyxminiproblems-problemphyxmini0608-speedquantity, def:physics:phyx-mini-0608:phyxminiproblems-problemphyxmini0608-starshipaffiliation, def:physics:phyx-mini-0608:phyxminiproblems-problemphyxmini0608-ellipseaxisrole, def:physics:phyx-mini-0608:phyxminiproblems-problemphyxmini0608-inertialframelabel, def:physics:phyx-mini-0608:phyxminiproblems-problemphyxmini0608-starshipmarkingfigure}
  Independent physical quantities in the observation. The observed axis lengths and relative speed are fields rather than definitions from the contraction law or an answer choice
\end{definition}
\begin{proof}
  This is the declared model definition of Empire Marking Observation.
\end{proof}
\begin{definition}[speed Fraction Of Light]
  \label{def:physics:phyx-mini-0608:phyxminiproblems-problemphyxmini0608-speedfractionoflight}
  \lean{PhyXMiniProblems.ProblemPhyXMini0608.speedFractionOfLight}
  \uses{def:physics:phyx-mini-0608:phyxminiproblems-problemphyxmini0608-speedinmeterspersecond, def:physics:phyx-mini-0608:phyxminiproblems-problemphyxmini0608-vacuumspeedoflightinmeterspersecond, def:physics:phyx-mini-0608:phyxminiproblems-problemphyxmini0608-empiremarkingobservation}
  The dimensionless relative speed β = v / c, evaluated in SI units
\end{definition}
\begin{proof}
  This is the declared model definition of speed Fraction Of Light.
\end{proof}
\begin{definition}[lorentz Factor]
  \label{def:physics:phyx-mini-0608:phyxminiproblems-problemphyxmini0608-lorentzfactor}
  \lean{PhyXMiniProblems.ProblemPhyXMini0608.lorentzFactor}
  \uses{def:physics:phyx-mini-0608:phyxminiproblems-problemphyxmini0608-empiremarkingobservation, def:physics:phyx-mini-0608:phyxminiproblems-problemphyxmini0608-speedfractionoflight}
  Physlib's Lorentz factor for the observation's dimensionless speed
\end{definition}
\begin{proof}
  This is the declared model definition of lorentz Factor.
\end{proof}
\begin{definition}[Matches Empire Marking Scenario]
  \label{def:physics:phyx-mini-0608:phyxminiproblems-problemphyxmini0608-matchesempiremarkingscenario}
  \lean{PhyXMiniProblems.ProblemPhyXMini0608.MatchesEmpireMarkingScenario}
  \uses{def:physics:phyx-mini-0608:phyxminiproblems-problemphyxmini0608-empiremarkingobservation}
  Frame, affiliation, and motion-axis roles stated or required by the setup
\end{definition}
\begin{proof}
  This is the declared model definition of Matches Empire Marking Scenario.
\end{proof}
\begin{definition}[Matches Supplied Starship Figure]
  \label{def:physics:phyx-mini-0608:phyxminiproblems-problemphyxmini0608-matchessuppliedstarshipfigure}
  \lean{PhyXMiniProblems.ProblemPhyXMini0608.MatchesSuppliedStarshipFigure}
  \uses{def:physics:phyx-mini-0608:phyxminiproblems-problemphyxmini0608-starshipmarkingfigure}
  Qualitative geometry and literal labels visible in test\_image/608.png
\end{definition}
\begin{proof}
  This is the declared model definition of Matches Supplied Starship Figure.
\end{proof}
\begin{definition}[Matches Problem Ellipse Ratio]
  \label{def:physics:phyx-mini-0608:phyxminiproblems-problemphyxmini0608-matchesproblemellipseratio}
  \lean{PhyXMiniProblems.ProblemPhyXMini0608.MatchesProblemEllipseRatio}
  \uses{def:physics:phyx-mini-0608:phyxminiproblems-problemphyxmini0608-lengthreadout, def:physics:phyx-mini-0608:phyxminiproblems-problemphyxmini0608-empiremarkingobservation}
  The dimensionless a = 1.40 b datum from the prose, linked to the two dimensionful proper axis lengths in every common length unit
\end{definition}
\begin{proof}
  This is the declared model definition of Matches Problem Ellipse Ratio.
\end{proof}
\begin{definition}[Has Physical Empire Marking Parameters]
  \label{def:physics:phyx-mini-0608:phyxminiproblems-problemphyxmini0608-hasphysicalempiremarkingparameters}
  \lean{PhyXMiniProblems.ProblemPhyXMini0608.HasPhysicalEmpireMarkingParameters}
  \uses{def:physics:phyx-mini-0608:phyxminiproblems-problemphyxmini0608-lengthreadout, def:physics:phyx-mini-0608:phyxminiproblems-problemphyxmini0608-vacuumspeedoflightinmeterspersecond, def:physics:phyx-mini-0608:phyxminiproblems-problemphyxmini0608-empiremarkingobservation, def:physics:phyx-mini-0608:phyxminiproblems-problemphyxmini0608-speedfractionoflight}
  Positivity and the nonzero subluminal branch relevant to the question
\end{definition}
\begin{proof}
  This is the declared model definition of Has Physical Empire Marking Parameters.
\end{proof}
\begin{definition}[Appears As Federation Circle]
  \label{def:physics:phyx-mini-0608:phyxminiproblems-problemphyxmini0608-appearsasfederationcircle}
  \lean{PhyXMiniProblems.ProblemPhyXMini0608.AppearsAsFederationCircle}
  \uses{def:physics:phyx-mini-0608:phyxminiproblems-problemphyxmini0608-lengthreadout, def:physics:phyx-mini-0608:phyxminiproblems-problemphyxmini0608-empiremarkingobservation}
  The requested observational condition: the moving ellipse has equal apparent axis lengths and can therefore be confused with the Federation circle. This does not supply the speed needed to produce that appearance
\end{definition}
\begin{proof}
  This is the declared model definition of Appears As Federation Circle.
\end{proof}
\begin{definition}[Satisfies Relativistic Axis Contraction]
  \label{def:physics:phyx-mini-0608:phyxminiproblems-problemphyxmini0608-satisfiesrelativisticaxiscontraction}
  \lean{PhyXMiniProblems.ProblemPhyXMini0608.SatisfiesRelativisticAxisContraction}
  \uses{def:physics:phyx-mini-0608:phyxminiproblems-problemphyxmini0608-lengthreadout, def:physics:phyx-mini-0608:phyxminiproblems-problemphyxmini0608-empiremarkingobservation, def:physics:phyx-mini-0608:phyxminiproblems-problemphyxmini0608-lorentzfactor}
  The generic special-relativistic length law. The axis parallel to the motion contracts by 1 / γ(β) and the perpendicular axis is unchanged. Neither clause specializes β to this problem's numerical result
\end{definition}
\begin{proof}
  This is the declared model definition of Satisfies Relativistic Axis Contraction.
\end{proof}
\begin{definition}[Answer Choice]
  \label{def:physics:phyx-mini-0608:phyxminiproblems-problemphyxmini0608-answerchoice}
  \lean{PhyXMiniProblems.ProblemPhyXMini0608.AnswerChoice}
  Labels of the four speed choices in the source problem
\end{definition}
\begin{proof}
  This is the declared model definition of Answer Choice.
\end{proof}
\begin{definition}[displayed Speed In Meters Per Second]
  \label{def:physics:phyx-mini-0608:phyxminiproblems-problemphyxmini0608-displayedspeedinmeterspersecond}
  \lean{PhyXMiniProblems.ProblemPhyXMini0608.displayedSpeedInMetersPerSecond}
  \uses{def:physics:phyx-mini-0608:phyxminiproblems-problemphyxmini0608-answerchoice}
  Displayed speed in metres per second for each answer choice
\end{definition}
\begin{proof}
  This is the declared model definition of displayed Speed In Meters Per Second.
\end{proof}
\begin{definition}[recorded Dataset Answer]
  \label{def:physics:phyx-mini-0608:phyxminiproblems-problemphyxmini0608-recordeddatasetanswer}
  \lean{PhyXMiniProblems.ProblemPhyXMini0608.recordedDatasetAnswer}
  \uses{def:physics:phyx-mini-0608:phyxminiproblems-problemphyxmini0608-answerchoice}
  The answer label recorded by the source dataset
\end{definition}
\begin{proof}
  This is the declared model definition of recorded Dataset Answer.
\end{proof}
\begin{definition}[Matches Displayed Speed Choice]
  \label{def:physics:phyx-mini-0608:phyxminiproblems-problemphyxmini0608-matchesdisplayedspeedchoice}
  \lean{PhyXMiniProblems.ProblemPhyXMini0608.MatchesDisplayedSpeedChoice}
  \uses{def:physics:phyx-mini-0608:phyxminiproblems-problemphyxmini0608-speedinmeterspersecond, def:physics:phyx-mini-0608:phyxminiproblems-problemphyxmini0608-empiremarkingobservation, def:physics:phyx-mini-0608:phyxminiproblems-problemphyxmini0608-answerchoice, def:physics:phyx-mini-0608:phyxminiproblems-problemphyxmini0608-displayedspeedinmeterspersecond}
  Agreement with a speed displayed as 2.10 × 10⁸ m/s. The tolerance of 5 × 10⁵ m/s is half one unit in the final displayed digit, so this predicate expresses rounding to three significant figures
\end{definition}
\begin{proof}
  This is the declared model definition of Matches Displayed Speed Choice.
\end{proof}
% --- Archon declaration topology end ---
% --- Archon physics formalization source end ---
